\documentclass[11pt]{article}
\usepackage{amsmath,amssymb,amsthm}
\newtheorem{proposition}{Proposition}
\newtheorem{theorem}[proposition]{Theorem}
\begin{document}
\title{Signed Logarithmic Potentials in the Positive Denominator Chamber}
\author{Head-compression research audit}
\maketitle

\paragraph{Status and scope.}
This note gives a complete written construction and a small exact arithmetic
check. It is not formalized in Lean. It is not an upper or lower bound for
\texttt{worstCaseHeadComplexity}. Its purpose is to test a proposed route to
the global lower bound: positivity of denominators does not, by itself,
make the associated signed logarithmic master functions have only
$C^n$ critical points. The actual interpolation-fiber and strict-sign
counting questions remain unresolved.

For weights $\lambda_x\in\mathbb R$ indexed by $x\in\{0,1\}^n$, put
\[
 \Phi_\lambda(c)=\sum_x\lambda_x\log(1+c\cdot x),
 \qquad c\in(0,\infty)^n.
\]
All logarithms are natural. Every affine form appearing here has positive
constant coefficient and strictly positive slopes in the coordinates
where the corresponding $x_i$ is one. The parameter $c$ ranges inside
the strict positive-slope chamber of the original affine denominator.

\begin{proposition}[Elliptic identity]
Set $Q(c)=\operatorname{diag}(c)+cc^{\mathsf T}$. Then $Q(c)$ is positive
definite and
\[
 \operatorname{tr}\bigl(Q(c)\nabla^2\Phi_\lambda(c)\bigr)
       +c\cdot\nabla\Phi_\lambda(c)=0.
\]
Consequently every nondegenerate critical point is a saddle.
\end{proposition}
\begin{proof}
For a nonzero vector $v$, one has
$v^{\mathsf T}Q(c)v=\sum_i c_i v_i^2+(c\cdot v)^2>0$.
For $B_x=1+c\cdot x$,
\[
 \nabla\log B_x=x/B_x,\qquad
 \nabla^2\log B_x=-xx^{\mathsf T}/B_x^2.
\]
The Boolean identity $x_i^2=x_i$ gives
$x^{\mathsf T}Q(c)x=c\cdot x+(c\cdot x)^2=(c\cdot x)B_x$.
The asserted differential identity follows term by term. At a critical
point, a positive or negative definite Hessian would have a nonzero trace
against $Q(c)$ of the corresponding sign. A nonsingular symmetric Hessian
with zero such trace therefore has eigenvalues of both signs.
\end{proof}

\begin{theorem}[Many critical points in one positive chamber]
Let $d\geq2$ be even and let $k\geq1$. For $n=d+k$ there are real weights
$\lambda_x$ such that $\Phi_\lambda$ has at least
\[
                  (2^d-1)^k
\]
distinct nondegenerate critical points in $(0,\infty)^n$.
In particular there is an absolute $c>0$ such that, for every sufficiently
large $n$, some such potential has at least $2^{c n^2}$ nondegenerate
critical points in the positive chamber.
\end{theorem}
\begin{proof}
Write the parameters as $(b_1,\ldots,b_k,c_1,\ldots,c_d)$.
Choose $c^0_j=2^{j-1}$, so the numbers
\[
       \beta_y=1+c^0\cdot y,\qquad y\in\{0,1\}^d,
\]
are the $L=2^d$ distinct integers $1,\ldots,L$.
Let $P(t)=\prod_{r=1}^{L-1}(t-r)$ and
$Q_0(t)=\prod_y(t+\beta_y)$. Partial fractions give
\[
 \frac{P(t)}{Q_0(t)}=\sum_y\frac{a_y}{t+\beta_y},
 \qquad
 a_y=\frac{P(-\beta_y)}{Q_0'(-\beta_y)}.
\]
Define
\[
 f_i(b_i,c)=\sum_y a_y\log(1+b_i+c\cdot y).
\]
At $c=c^0$, its derivative in $b_i$ is $P(b_i)/Q_0(b_i)$.
It has the simple positive roots $1,\ldots,L-1$; the second derivative at
root $r$ is $P'(r)/Q_0(r)\neq0$.

We next pin the address parameters near $c^0$ by a nondegenerate saddle.
Partition their indices into pairs. For a pair $(u,v)$ set
$U=1+c^0_u$, $V=1+c^0_v$, and $B=1+c^0_u+c^0_v$, and put
\[
 \psi_{u,v}(c_u,c_v)=\log(1+c_u+c_v)
       -\frac{U}{B}\log(1+c_u)-\frac{V}{B}\log(1+c_v).
\]
Its gradient vanishes at $(c^0_u,c^0_v)$. Its Hessian there is
\[
 \frac1{B^2}
 \begin{pmatrix}c^0_v/U&-1\\-1&c^0_u/V\end{pmatrix},
 \qquad
 \det\nabla^2\psi_{u,v}=-\frac1{B^3UV}\neq0.
\]
The sum $\psi$ of these pair potentials therefore has a nondegenerate
critical point at $c^0$.

For $\varepsilon>0$ consider
\[
 \Phi_\varepsilon(b,c)=\varepsilon^{-1}\psi(c)
                       +\sum_{i=1}^k f_i(b_i,c).
\]
This is a potential of the required form. The $f_i$ terms use cube
vertices with signal block $e_i$; the $\psi$ terms use vertices with
signal block zero. Repeated terms may simply have their weights added.

Consider the smooth system
\[
 G_\varepsilon(b,c)=
 \left(\partial_{b_1}f_1,\ldots,\partial_{b_k}f_k,
       \nabla\psi(c)+\varepsilon\sum_i\nabla_c f_i(b_i,c)\right)=0.
\]
For $\varepsilon>0$ this is exactly the critical-point system for
$\Phi_\varepsilon$, after multiplying its last $d$ equations by
$\varepsilon$. At $\varepsilon=0$, every point
\[
       (r_1,\ldots,r_k,c^0),\qquad r_i\in\{1,\ldots,L-1\},
\]
is a solution. The Jacobian in $(b,c)$ is block upper triangular.
Its diagonal blocks are the diagonal matrix with entries
$P'(r_i)/Q_0(r_i)$ and the invertible Hessian $\nabla^2\psi(c^0)$.
Thus every such solution is nonsingular.

The ordinary real implicit function theorem gives a nearby nonsingular
solution for all sufficiently small positive $\varepsilon$. Choose
disjoint neighborhoods of the finitely many initial points, all inside
the positive orthant, and then one common sufficiently small
$\varepsilon>0$. This gives $(L-1)^k$ distinct solutions. Nonsingularity
of $G_\varepsilon$ is equivalent to nonsingularity of the Hessian of
$\Phi_\varepsilon$, because their Jacobians differ by an invertible
diagonal scaling of the last $d$ rows.

Finally take the largest even $d\leq n/2$ and $k=n-d$. For sufficiently
large $n$, both are at least $n/3$, and $2^d-1\geq2^{d-1}$.
The count is at least $2^{(d-1)k}\geq2^{n^2/12}$ eventually.
\end{proof}

\paragraph{Small exact check.}
The script \texttt{reverse-alternative-master-check.py} verifies the
$d=2$, $k=2$ limiting system using rational arithmetic. It takes
$c^0=(1,2)$, $\beta=(1,2,3,4)$, and residues
$a=(-4,30,-60,35)$. The three signal second derivatives are
$1/60,-1/360,1/420$. The address Hessian determinant is $-1/384$.
The script checks all nine nonsingular limiting points. It does not
replace the implicit-function argument or assert a numerical universal
compiler.

\paragraph{Relation to head counting.}
For the one-head evaluation vector $x\mapsto A(x)/B(x)$, a linear
functional with coefficients $\lambda_x$ has derivative in the numerator
parameters $\sum_x\lambda_x(1,x)/B(x)$. Its derivative in denominator
parameters is $-\sum_x\lambda_x A(x)(1,x)/B(x)^2$.
In homogeneous notation these are the gradient of the logarithmic
potential and its Hessian applied to the numerator coefficient vector.
At a nonconstant normal contact, the Hessian must have a kernel beyond
the scaling direction. The theorem above concerns nondegenerate
critical points after normalization. It therefore does not directly
construct many nonconstant normal contacts or interpolation fibers.
It only refutes a bound on all positive-chamber signed master-function
critical points of the form $C^n$. No claim about the desired
$H!C^{Hn}$ interpolation or sign-region bound follows in either direction.

\paragraph{The full Boolean differential identity.}
The trace identity is only one consequence of the stronger system
\[
 \partial_{ii}\Phi+\sum_jc_j\partial_{ij}\Phi+\partial_i\Phi=0
 \qquad(1\leq i\leq n).
\]
It follows termwise from $x_i^2=x_i$, since the three terms for
$\log B_x$ add to $-x_i/B_x^2-x_i(c\cdot x)/B_x^2+x_i/B_x=0$.
At a critical point the Hessian $H$ therefore satisfies
$\operatorname{diag}(H)+Hc=0$. This is stronger than the weighted trace
condition, but does not force every degenerate critical point to have
large corank.

\begin{proposition}[An ordinary nonconstant normal contact]
There is a signed master potential on the three-dimensional Boolean cube
with a critical point at $c=(1,1,1)$ whose Hessian has corank one and whose
third derivative along its kernel is nonzero. The associated nonconstant
one-head normal contact is a nondegenerate critical point after fixing
the constant and projective scaling directions.
\end{proposition}
\begin{proof}
In vertex order $000,100,010,001,110,101,011,111$, take
\[
 \lambda=(3,12,-12,-12,-9,-9,27,0),\qquad v=(2,1,1).
\]
Direct rational arithmetic gives $\sum_x\lambda_x=0$,
$\nabla\Phi(1,1,1)=0$, and
\[
 H=\nabla^2\Phi(1,1,1)=
 \begin{pmatrix}-1&1&1\\1&1&-3\\1&-3&1\end{pmatrix},
 \qquad Hv=0,\qquad
 D^3\Phi(1,1,1)[v,v,v]=-2.
\]
The two nonzero eigenvalues of $H$ are $4$ and $-3$.
The identity $\sum\lambda_x=0$ also ensures that the homogeneous
gradient in all affine denominator coefficients vanishes.
Consequently the affine-ratio score
$A(x)/B(x)$, with $A(x)=2x_1+x_2+x_3$ and
$B(x)=1+x_1+x_2+x_3$, is a nonconstant normal contact for the
functional $G\mapsto\sum_x\lambda_xG(x)$. The denominator has
strictly positive slopes and positive constant coefficient.

For a direct regularity check remove the irrelevant constant output and
fix the coefficient of $x_1$ in $A$ to two. Use local parameters
$(a_2,a_3,c_1,c_2,c_3)$ for
\[
 L(a,c)=\sum_x\lambda_x
       \frac{2x_1+a_2x_2+a_3x_3}{1+c\cdot x}.
\]
The point $(1,1,1,1,1)$ is critical. Its Hessian has upper-left
$2$ by $2$ block zero, upper-right block consisting of the last two
rows of $H$, and lower-right block
\[
 T=D^3\Phi[v]=
 \begin{pmatrix}2&-2&-2\\-2&-1&4\\-2&4&-1\end{pmatrix}.
\]
The resulting $5$ by $5$ determinant is $-32$, so the critical point is
nondegenerate in this chart. Adding one to the score gives the numerator
$A+B$, which is strictly positive on the entire cube. Since
$\sum\lambda_x=0$, this addition changes neither the contact nor the
regularity calculation. All equalities are checked by the accompanying
exact-arithmetic script.
\end{proof}

\begin{lemma}[Regularity of a normalized contact]
Suppose $\nabla\Phi(c)=0$, the Hessian $H$ has kernel
$\mathbb R a$, and $D^3\Phi(c)[a,a,a]\neq0$. Choose an index $\ell$
with $a_\ell\neq0$. The function
\[
 (u,b)\longmapsto u\cdot\nabla\Phi(b),\qquad u_\ell=a_\ell,
\]
has a nondegenerate critical point at $(a,c)$.
\end{lemma}
\begin{proof}
The first derivatives vanish because $\nabla\Phi(c)=0$ and $Ha=0$.
Let $U=\{p:p_\ell=0\}$. The Hessian on $U\times\mathbb R^n$ has
block form
\[
 \begin{pmatrix}0&H|_{U\times\mathbb R^n}\\
 H|_{\mathbb R^n\times U}&D^3\Phi(c)[a]\end{pmatrix}.
\]
If $(p,q)$ belongs to its kernel, the upper equations say that
$Hq$ is a multiple of the coordinate vector $e_\ell$.
Since $a^{\mathsf T}Hq=0$ and $a_\ell\neq0$, that multiple is zero.
Thus $q=ta$. Pairing the lower equations with $a$ gives
$tD^3\Phi(c)[a,a,a]=0$, hence $t=0$. The lower equations now give
$Hp=0$, so $p$ is a multiple of $a$. Since $p_\ell=0$, it is zero.
The Hessian is therefore nonsingular.
\end{proof}

\begin{theorem}[Many normalized nonconstant contacts]
For every sufficiently large dimension $N$ there is a signed linear
functional on full-cube evaluation vectors having at least
$2^{cN^2}$ distinct nonconstant one-head normal contacts, for an absolute
$c>0$. All denominators have positive constant and strictly positive
slopes. The scores can all be chosen strictly positive on the entire
cube. Each contact is isolated and nondegenerate after quotienting out
the constant-output and output-scaling directions.
\end{theorem}
\begin{proof}
Apply the many-critical-points theorem in dimension $n=N-3$ and denote
the resulting potential by $\Phi_{\mathrm{many}}(b)$.
Let $\Phi_{\mathrm{fold}}(d)$ be the three-variable potential in the
preceding proposition. The sum
\[
 \Phi(b,d)=\Phi_{\mathrm{many}}(b)+\Phi_{\mathrm{fold}}(d)
\]
is again a signed Boolean-cube master potential: embed the first set of
weights on vertices whose last three coordinates vanish and the second
set on vertices whose first $n$ coordinates vanish. Their origin weights
are added. Adjust the weight at the origin so that the total weight is
zero; this changes neither potential nor any of its derivatives.

Every nondegenerate critical point $b^*$ of
$\Phi_{\mathrm{many}}$, paired with $d=(1,1,1)$, is critical for
$\Phi$. Its Hessian is block diagonal, with the invertible first block
and the corank-one three-variable block. Its kernel is therefore spanned
by $a=(0,\ldots,0,2,1,1)$, and its third derivative along $a$ is $-2$.
The regularity lemma applies at every such point.

For the associated cube functional $\lambda$, consider the head
\[
 G_{b^*}(x,y)=
 \frac{2y_1+y_2+y_3}{1+b^*\cdot x+y_1+y_2+y_3}.
\]
Derivative in every numerator slope is a component of
$\nabla\Phi=0$, and derivative in denominator slopes is a component
of $Ha=0$. Derivative in the numerator constant is also zero, because
\[
 \sum_z\frac{\lambda_z}{1+c\cdot z}
 =\sum_z\lambda_z-c\cdot\nabla\Phi(c)=0.
\]
The homogeneous denominator-scaling derivative then vanishes as well.
Thus this is a normal contact for the full affine-ratio family.
The chart used in the lemma fixes numerator constant zero and the first
nonzero numerator slope to two. These remove precisely the locally
irrelevant constant-output and output-scaling directions. The lemma
shows that the contact is isolated and nondegenerate in this chart.

Different $b^*$ give different cube scores: at $y=(1,0,0)$ and
$x=e_i$, the value is $2/(2+b^*_i)$, which determines $b^*_i$.
Adding one makes every score strictly positive. The corresponding
numerator remains affine, and the zero total weight means this addition
does not affect any contact assertion. There are at least
$2^{n^2/12}$ such contacts for all sufficiently large $n$, which is
$2^{cN^2}$ for a fixed positive $c$ and all sufficiently large $N$.
\end{proof}

\paragraph{Scope of the contact theorem.}
This refutes a uniform $C^N$ bound on the number of normalized
nonconstant contacts of an arbitrary signed full-cube functional in
the positive denominator chamber. It does not refute the proposed
$H!C^{HN}$ interpolation or strict-sign count. The functional constructed
here is special. Under a perturbation of the functional, the
nondegenerate critical points in the normalized chart persist, but
their critical values need not remain zero. Consequently these
perturbations need not remain normal contacts of the unnormalized cone
of head scores. No generic conic-normal-fiber statement is asserted.
All the examples above can lie in the same strictly positive output
orthant, so a large contact count alone also gives no strict-sign
shattering theorem. The new theorem is a written proof, not a Lean
formalization, and neither main asymptotic target is settled by it.
\end{document}
